\chapter{The simulator}
\label{ch:simulator}

The simulator runs Modbus devices inside Modbench. Use it to try a map, an
alarm or another program (a SCADA, a PLC, a script) without the hardware:
each device answers like the real one, with the values of its map, and can
make values move on its own.

Open it from \nav{Tools\then Simulator}.

\screenshot[9cm]{simulator}{The simulator with two devices on a serial port.}
  {The simulator running on Modbus TCP: two device tabs (Unit 1 and Unit 2),
   the map picker, the registers view with values, and the recent requests at
   the bottom.}

\section{Where it listens}

\begin{description}[style=nextline, leftmargin=0pt, font=\intermedium]
  \item[Interface] Modbus TCP or UDP, RTU over TCP or UDP, or a serial port
    with RTU or ASCII. Over the network, set the \field{Port} (5020 by
    default, so it does not need administrator rights); \field{Other
    computers} also accepts clients from the network, not only from this
    computer. On a serial port, set the line settings as the clients use them.
  \item[\uiaccent{Start} and \ui{Stop}] Start listens; the status shows where
    (\emph{Listening on 127.0.0.1:5020}).
  \item[\ui{Connect Main to it}] Fills step 1 of the Main page with the
    simulator's address and the shown device's unit ID, and goes to Main.
\end{description}

\begin{note}[Serial]
Clients connect to the other end of the line. To try it on one computer, use
a virtual null-modem pair (com0com on Windows, \val{socat} on Linux and macOS)
and open the other port on the Main page.
\end{note}

\section{Devices}

The simulator starts with one device. Several devices can share one address,
each on its own unit ID, like devices behind a gateway or on one RS-485 line;
other unit IDs get no answer.

\begin{description}[style=nextline, leftmargin=0pt, font=\intermedium]
  \item[\ui{Add device}] Adds a device on the lowest free unit ID. The tabs
    above switch between devices.
  \item[\field{Unit ID}] The unit ID the device answers. Empty answers every
    unit ID that no other device has. Unit ID \val{0} is a broadcast: every
    device takes the write and none answers.
  \item[\field{Map}] Shapes the device after a map: one open on the Main page,
    one from the library, or \nav{Open a map file…}. \emph{Empty device}
    gives a memory where every address exists and holds zero.
  \item[\field{Only the map's addresses}] Addresses the map does not have
    answer exception 2, like the real device. Off, every address answers.
  \item[\ui{Reset values}] Puts back the map's reference values.
  \item[\ui{Remove device}] Removes the device shown; its unit ID gets no
    answer.
\end{description}

\section{Values}

The view switch shows the device's memory either as \tab{Registers}, the map's
registers by name with their values in engineering units, or as a raw
\tab{Grid} of addresses, in any area. Click a value to change it: clients
read the new value from the next request.

Writes from clients show up immediately, and the recent requests are listed
at the bottom, reads and writes apart, so a stream of reads does not push the
writes out of sight.

\section{Generators}

Generators make values change on their own, every so many milliseconds (at
least 50). \ui{Add generator} adds one; it drives a map register, in its own
units, or a raw address.

\begin{center}
\begin{tabularx}{\linewidth}{l L}
  \toprule
  \thead{Kind} & \thead{What it does} \\
  \midrule
  Increment & Adds a step each time, from a minimum to a maximum, then starts
    again: counters, energy. \\
  Sine & Swings around an offset with an amplitude and a period: voltage,
    temperature. \\
  Random & A new value between a minimum and a maximum each time: noise. \\
  Copy & Takes the value of another address: a status input that follows a
    command coil. \\
  \bottomrule
\end{tabularx}
\end{center}

\screenshot[5cm]{generators}{Generators on a meter's registers.}
  {The generators list: sines on the voltages, currents and powers, random
   values on the power factors and an increment on the energy.}

A generator that cannot run (its register is not in the map, or its address
does not exist) says why. Generators keep running across \ui{Stop} and
\uiaccent{Start}.

\section{Its own window}

\ui{Pop out} moves the simulator into a window of its own, to watch it next to
the Main page. \ui{Back to the main window} brings it back to the Tools page.

\section{Saving a setup}
\label{sec:sim-setups}

A test bench with several devices, maps and generators takes a while to set
up. \ui{Save setup…} keeps it in a \file{.mbsim} file:

\begin{itemize}
  \item the interface and its settings;
  \item every device: unit ID, map, \emph{Only the map's addresses}, and
        generators.
\end{itemize}

\ui{Open setup…}, or a double-click on the file, puts it back. Each map is read
again from its file when that is still there, so updates to the library are
picked up; otherwise the copy kept in the setup is used. While the simulator
runs, opening a setup replaces the devices but keeps the interface.

\begin{note}
Values changed by hand are not saved: devices start from their maps'
reference values.
\end{note}
